\begin{afterword}
\summaryhead

The post follows the previous day's ``Bystander Apathy.'' It suggests that the bystander effect may come from coordination instincts suited to groups of people who knew each other, and notes that the effect diminishes when subjects know each other. It then says that people seem to react poorly to problems met on the Internet, and lists eight possible reasons: online groups are strangers of unknown number, without eye contact or real-time talk, not dependent on each other, anonymous, surrounded by other people doing nothing, and hearing no voiced plea.

The author has no solution, and says that online activism tools tend to be petitions. The post offers some ideas: a video of someone asking for help, names of the first helpers, video thank-yous, referral codes that show how many people a helper brought in, and bounties for verifiable subtasks. It ends by handing readers the open problem of helping strangers form effective task forces online, noting Clay Shirky's estimate of how much spare time Americans spend watching television.

\responsehead

The post is a call to work on a problem, and it is candid about that. Its explanations come with ``may'' and ``Perhaps,'' its ideas are what ``come to mind,'' and it ends by handing the problem on. What can be asked of it is whether its framing fits the problem and how much of it was new.

The premise, that people ``seem to have a hard time reacting constructively to problems encountered over the Internet,'' is stated as obvious, without an example. The list of eight reasons is a set of plausible hypotheses, not tested in the post. One had been tested nine years earlier. Markey (2000) found that in chat groups help took longer as more people were present, and that the effect was virtually eliminated when the request named one person. That supports the post, and it is Cialdini's advice from the day before, applied online.

The remark that online activism tools ``tend to be'' petitions is hedged. At least one tool had been built for exactly this problem. PledgeBank, running since 2005, let anyone pledge to act if enough others would, in order to get people past ``the fear of acting alone.''

A closer model than the bystander research is Mancur Olson's. The bystander research the post builds on concerns emergencies, where people must decide quickly whether something is wrong and whether it is their job. The tasks the post wants done (donations, open-source code, bounties) are continuing contributions to a shared good, the case Olson analyzed in 1965 as a problem of incentives in large groups. Several of the post's proposals, public credit and thanks for those who contribute, are what Olson called selective incentives: rewards given only to those who contribute. The post's framing in terms of instincts and Olson's in terms of incentives lead here to similar remedies.

The Shirky figure is reported correctly, as a comparison with one year of American television.

In short: a hedged call to build tools against collective inaction online, whose hypotheses and proposals had already been partly tested or tried: in a 2000 study of chat groups, in PledgeBank, and in Olson's account of selective incentives.
\end{afterword}
